\begin{abstract}
We disprove the unrestricted four-dimensional disk-embedding conjecture.
We construct immersed disks in a compact oriented smooth four-manifold with
framed algebraic dual spheres satisfying the usual equivariant intersection
and reduced self-intersection conditions, but with no pairwise disjoint locally
flat replacements that preserve the boundary maps and induced normal framings.
The obstruction holds even when the replacement disks' relative homotopy classes
are not prescribed.
\end{abstract}

\section{Introduction}\label{sec:introduction}

The disk-embedding problem asks when immersed disks in a four-manifold can
be replaced by disjoint locally flat embedded disks with the same boundary
data. Its role in four-dimensional topological surgery is to turn
algebraic cancellation data into embedded disks.
In dimension four, the auxiliary disks used to carry out cancellations
can themselves intersect the original surfaces and one another.
Dual spheres provide additional control. An algebraic dual records one
intersection with its designated disk and zero with the others, counted
with signs and fundamental-group labels; a geometric dual realizes these
counts by meeting its designated disk once transversely and missing the
other disks.

Casson's construction replaces the finite task of removing intersections
by an infinite tower of immersed disks. Freedman's recognition theorem
identifies the resulting Casson handle, relative to its attaching region,
with the standard open topological two-handle. Together these results
established the simply connected foundation of disk embedding
\cite[Introduction and Theorem~1.1]{Freedman1982}.
Freedman and Quinn developed the capped-grope version of this method and
organized the disk theorem and its surgery and cobordism applications
\cite[Theorem~5.1A and Chapter~11]{FQ1990}.
The group restriction enters through the double-point loops of the caps:
their null-homotopies supply further stages of the construction.
The class of \emph{good groups} is defined by the availability of the
required \(\pi_1\)-null disk operations
\cite[Definition~3.2]{PRT2025}.
Freedman and Teichner proved that groups of subexponential growth are good
\cite[Theorem~0.1]{FT1995}. Krushkal and Quinn supplied another proof
\cite[Theorem, p.~408]{KQ2000}. Their paper also contains
Freedman and Teichner's correction of the earlier height-raising
argument \cite[Appendix, pp.~424--430]{KQ2000}.

The modern formulation of Powell, Ray, and Teichner supplies framed
geometric dual spheres as well as embedded disks, and preserves the
homotopy classes of the input dual spheres
\cite[Theorem~A]{PRT2025}. Their modification supplies the
geometric-dual step missing from the earlier argument; the disk
conclusion without these output spheres was already proved
\cite[Remark~1.3]{PRT2025}.
For generically immersed input disks, their theorem preserves the induced
boundary normal framings without prescribing relative homotopy classes
for the replacement disks.
The \emph{unrestricted disk-embedding conjecture} considered here is the
extension of that theorem obtained by removing only the good-group
hypothesis. It would in particular give framed disk replacements for
compact oriented smooth input manifolds, even after the output-dual
requirement is discarded. We resolve this geometric conjecture negatively
by obstructing that disk conclusion.

We use the following intersection conventions in stating the result.
After choosing basepoints and whiskers, \(\lambda\) denotes the
equivariant intersection pairing with values in \(\mathbb Z[\pi_1M]\).
For an oriented four-manifold \(M\), the reduced self-intersection invariant
\(\widetilde\mu\) takes values in the additive quotient
\[
 \mathbb Z[\pi_1M]\big/
 \big\langle h-h^{-1}\ (h\in\pi_1M),\ \mathbb Z\cdot1\big\rangle .
\]
A framed immersed sphere has a trivialized normal bundle. A boundary
normal framing is the framing induced by the normal bundle of the
specified immersed disk.

\begin{theorem}\label{thm:main}
There exist a compact connected oriented smooth four-manifold \(M\) with boundary,
an integer \(k\geq1\), generically immersed proper disks
\[
 f_1,\ldots,f_k:(D^2,S^1)\longrightarrow(M,\partial M)
\]
with disjoint embedded boundary circles, and generically immersed framed spheres
\(g_1,\ldots,g_k:S^2\longrightarrow M\), such that
\[
 \lambda(f_i,g_j)=\delta_{ij},\qquad
 \lambda(g_i,g_j)=0\quad\text{for all }i,j,\qquad
 \widetilde\mu(g_i)=0\quad\text{for all }i,
\]
but there are no pairwise disjoint locally flat embedded disks in \(M\)
with the same boundary maps and the boundary normal framings induced by the
\(f_i\).
\end{theorem}

The diagonal intersection equations are part of the statement.
The conclusion of Theorem~\ref{thm:main} requires no homotopy constraint on a
hypothetical replacement disk.
The markings used in the proof specify based curves and their paths in
the fixed construction; they impose no further condition on hypothetical
replacement disks. The result concerns geometric disk replacement.
The relation to independently formulated surgery assertions is discussed
after the proof in Section~\ref{sec:conclusion}.

\subsection{Idea of the proof}

We first construct a fixed immersed disk problem in the exterior of a
finite collection of capped gropes. A capped grope here is a surface tree
with disks attached at its terminal curves; its boundary records an
iterated commutator in the terminal curves. The construction installs
finitely many such commutator relations, with independently chosen based
conjugations, in a smooth plumbing model. It also supplies the algebraic
dual spheres required in Theorem~\ref{thm:main}. These data are fixed
before any hypothetical embedded replacements are chosen.

Suppose this particular disk problem has replacements with the required
boundary maps and framings. They allow the plumbing model to be cut into
two vertex pieces \(A,B\), joined by compact three-manifolds
\(Y_1,\ldots,Y_{2m}\) with torus boundary, for an integer \(m\geq5\)
chosen in the construction. The cuts are smooth and miss
the locally flat disks. The installed boundary words are now relations
in the vertex groups, and restriction in cohomology gives complementary
linear data on the two sides. Proposition~\ref{prop:geometric-reduction}
records this marked geometry; Lemma~\ref{lem:exterior-duals} verifies
the equivariant disk-and-sphere input.

For each edge \(Y_i\), we study based flat \(SL_2\)-connections with one
torus holonomy fixed. Basing records holonomy as a matrix, with no quotient
by simultaneous conjugation. A \emph{potential} is a function in finitely many
formal coordinates whose gradient equations describe these flat
connections. The \emph{central potential} fixes this torus
holonomy to the identity. Proposition~\ref{prop:edge-potential} constructs it from
the Chern--Simons functional, including a first-order variation of the
fixed boundary holonomy. Theorem~\ref{thm:edge-algebraicity} then replaces
the central potential by an algebraic function through a formal
coordinate change. The comparison preserves based holonomy on the full
critical scheme: its coordinate ring is the power-series ring modulo
the actual gradient ideal, with its nilpotent elements retained.
The geometric reduction selects \(m\) edges for the first-order
boundary variations; we call these edges active.

The two vertex pieces determine formal graphs in the product of the edge
coordinate spaces. Their central tangent spaces are complementary, and
suitable sums of the edge potentials vanish identically on the graphs.
The boundary relations imply more than vanishing at common critical
points: products associated to every three distinct active edges belong
to the restricted gradient ideals. On one side these products are
\(t_it_jt_k\), where each \(t_i^2=0\) and mixed products remain in the
parameter ring. On the other they are products of three algebraic
holonomy functions. Proposition~\ref{prop:vertex-data} establishes the
exact vanishing and ideal memberships. The graph parametrizations may
mix all edge coordinates, while each fixed potential depends only on
its own edge block.

Theorem~\ref{thm:tensor-obstruction} rules out precisely these algebraic
data when at least five edges are active. Artin approximation and finite
specialization first reduce them to large positive characteristic.
Truncating each internal coordinate by its characteristic power makes
the associated Koszul complexes finite. The two transverse graphs
determine classes with nonzero pairing. The target triple relations,
together with ordinary exterior differentiation, decompose one class
into tensors admitting three independent first-order parameter lifts.
The source triple relations force each resulting pairing to vanish,
giving the contradiction.

Three features are useful beyond the final contradiction. First, the geometry
imposes a finite family of independently conjugated laws before the cuts are
chosen. Second, the algebraicity argument compares full relative closed shifted
forms and preserves the marked holonomies; an algebraic model of the reduced
representation locus would be insufficient. Third, the tensor argument uses
actual gradient ideals over nonreduced parameter rings and makes no
isolated-critical-point or regular-sequence assumption.

\subsection{The geometric and algebraic antecedents}

Capped-grope contraction and transverse grope constructions provide the
local geometric operations \cite[\S\S1.2 and 1.4]{KQ2000}.
Their use to relate surface intersections and commutator words also
appears in the work of Freedman and Krushkal. Using the repeated-variable
2-Engel relation, they constructed homotopy solutions to the A--B slice
problem and new universal surgery models
\cite[Theorems~1 and~2]{FK2016}; their later work established universality
for half-\(\pi_1\)-null surgery models, with a capped grope on one side
and a \(\pi_1\)-null sphere on the other
\cite[Theorem~1 and Section~3]{FK2020}.
These results concern particular surgery or slicing formulations, with
their own embedding and complement conditions. They explain the value
of commutator calculus in this setting without supplying the marked
geometric reduction used here.

Our relations use three distinct group inputs and become trivial when
any one is the identity. Their logarithmic expansion begins with the
trilinear bracket \([Z_1,[Z_2,Z_3]]\), where \(Z_a\) is the logarithm
of the \(a\)-th input. The deletion property is the Brunnian property defined in
Section~\ref{sec:geometry}. We install a finite family of such words
with independently chosen conjugations, before knowing the cut pieces.
The later argument uses those independent choices to recover every
triple coefficient. It neither imposes a universal Engel relation nor
requires the vertex groups to be nilpotent. The simultaneous placement,
framing, and algebraic-dual calculations needed for this finite family
are given in Section~\ref{sec:geometry}.

The deformation-theoretic part has a different origin. Goldman and
Millson describe Artin-algebra families of flat connections by differential
graded Lie algebras and identify their based holonomies with representation
deformations \cite[Corollary~6.4 and Proposition~6.6]{GM1988}.
Their quadraticity results use additional formality hypotheses in the
compact K\"ahler setting. The edge manifolds here are arbitrary compact
three-manifolds, so higher deformation terms remain.
Chern and Simons constructed the classical transgression form whose
exterior derivative is the corresponding curvature characteristic form
\cite[Section~3, especially Proposition~3.2]{ChernSimons1974}.
Section~\ref{sec:edge-germs} uses this form to construct the potential,
computes its boundary adjustment, and proves equality of its full
curvature and gradient ideals.

The algebraic model comes from shifted symplectic geometry. The Betti
transgression theorem of Pantev, To\"en, Vaqui\'e, and Vezzosi produces
shifted symplectic forms on local-system stacks; their Lagrangian
intersection theorem lowers the shift by one
\cite[Theorems~2.5 and~2.9]{PTVV2013}.
Calaque's relative orientation construction supplies the Lagrangian
restriction map associated to an oriented manifold with boundary
\cite[Theorem~2.9 and \S3.1]{Calaque2015}.
After a punctured filling of the edge, these results put us in the
setting of Brav, Bussi, and Joyce's algebraic Darboux theorem in
shift \(-1\), which supplies an algebraic critical chart
\cite[Example~5.15 and Theorem~5.18(i)]{BBJ2019}.

To use that chart, we must identify its full relative closed form with
the form underlying the particular Chern--Simons potential already
constructed. Proposition~\ref{prop:formal-trace-comparison}, proved in
Appendix~\ref{app:formal-trace-comparison}, carries out this comparison
on derived coefficients and source cochain models. Its inputs include
Getzler's simplicial Maurer--Cartan comparison for nilpotent Lie algebras
\cite[Theorem~5.4 and Corollary~5.11]{Getzler2009} and the simplicial
de Rham contraction of Dupont \cite{Dupont1976}.
Section~\ref{sec:algebraicity} then turns the closed-form comparison
into a coordinate change fixing the critical scheme and its based
holonomies. Only after this step, and the polynomial replacement of the
holonomy logarithms, do we apply ordinary Artin approximation
\cite[Theorem~1.10]{Artin1969}. The approximation is applied to the
final graph equations with their fixed, separately algebraic edge
functions.

\subsection{Conventions and organization}

We assume familiarity with basic four-manifold topology, cohomology with
local coefficients, and formal deformation theory. The algebraicity
argument additionally uses derived mapping stacks and shifted symplectic
forms; its required comparison is stated in the body and proved in full
in Appendix~\ref{app:formal-trace-comparison}.

All manifolds and maps in the initial model are smooth, with corners rounded
when necessary. Hypothetical output disks are only locally flat. The subsequent
cuts are smooth hypersurfaces in the original smooth manifold and miss those
disks. Thus the vertex and edge calculations use ordinary smooth forms and
Stokes' theorem.

Unless a field is specified, cohomology has characteristic-zero coefficients.
We use \(\g=\mathfrak{sl}_2(\mathbb C)\) and the pairing
\(\kappa(X,Y)=\tr(XY)\). Connections and coordinate changes are formal power
series, coefficient by coefficient. A statement about an ideal is a statement
in the complete coordinate ring, including its nilpotents.
In Sections~\ref{sec:edge-germs}--\ref{sec:vertices}, \(d\) denotes the
differential on source forms and \(\delta\) differentiation in formal
coordinates. In the tensor argument of Section~\ref{sec:tensor}, \(d\) is
the coordinate de Rham differential and \(\delta\) its induced operator on
Koszul homology. We specify total signs when using both differentials.

Section~\ref{sec:geometry} proves the geometric reduction.
Section~\ref{sec:tensor} states and proves the finite obstruction before
we construct its inputs, making the target algebraic interface explicit.
Sections~\ref{sec:edge-germs} and \ref{sec:algebraicity}
construct the edge germs and their algebraic models.
Section~\ref{sec:vertices} constructs the vertex graphs and checks every
hypothesis of the tensor obstruction. Section~\ref{sec:conclusion} completes
the proof of Theorem~\ref{thm:main}.
